% 东方FX天则：规则书大纲
% 2026-10-16：按参与过程讲解；先说操作与结果，再解释术语。
% 保留既有规则与算例，不将设计参数写成已经确认的线上配置。

% 封面：金色金圆券与标注月虹市场的蓝色空白卡牌对照；引语落款“某粉毛股民妖怪”。
% 封面底部沿用原 rulebook 的 ISSUED BY / DOCUMENT 双栏样式。
\section*{FX，简单！}
\begin{itemize}
  \item 半页左右图解：顶部一行交代 500 本金和初价 5。
  \item 两张大卡片用 5 买 6 卖、5 卖 4 买的箭头例子介绍做多与做空，注明暂不计费用。
  \item 下排三张图例：玩家买入与央行注入金圆券分别推高价格；保证金率用“剩余本钱／交易规模”直观示意的大分数表示，并用全称提示强制平仓条件；限价单示意价格触线后自动尝试下单。三张卡片采用短标题，不附底部小字，图形、条件和结果按固定位置对齐。
  \item 底部简述消费计分、停盘安排及虚拟资产性质，后接目录。
\end{itemize}

\section{这场活动怎么玩}
\begin{itemize}
  \item 用买入、涨价、卖出说明基本玩法，交代交易池接单。
  \item 介绍交易与消费的关系，初始本金也能消费计分。
  \item 第一笔交易的操作流程，以及预热重置安排。
  \item 入门说明之后保留完整故事背景。
\end{itemize}

\section{买入和卖出}
\begin{itemize}
  \item 从 500 本金买入 20 张开始，再解释持仓、做多、平仓。
  \item 保留上涨和下跌后卖出的现金对照算例。
  \item 紧接着说明参考价与整笔均价的区别、订单对价格的影响、手续费。
  \item 保留不同订单大小的报价对照，以及持仓均价和预计平仓结果的用途。
  \item 多头盈亏未扣金圆券借款利息；部分卖出须看该笔订单报价。
\end{itemize}

\section{为什么价格会变}
\begin{itemize}
  \item 先讲玩家向交易池交入什么、取走什么，再解释买入推价和卖出压价。
  \item 统一使用价格；说明汇率在本活动中含义相同。
  \item 说明池子储备大小的影响，定价公式放附录。
  \item 比赛结果按事先公布的条件和额度触发投放；保留两种投放方向的算例。
  \item 投放不会直接给玩家发钱，也不会自动平仓或增加消费成绩。
\end{itemize}

\section{做空}
\begin{itemize}
  \item 先说明借卡卖出、跌价买回，以及涨价亏损和借卡利息。
  \item 从 500 本金借 20 张卖出讲起，保留涨跌后的现金对照算例；卖空所得不能直接当利润。
  \item 说明开空条件，以及只借卡牌也可能被强制平仓、亏损可能超过本金。
  \item 集中讲回补：含息数量、现金要求、部分回补及剩余欠款。
  \item 预计盈亏使用剩余空头对应的卖出所得，包含卡牌利息与交易成本，金圆券借款利息另看。
  \item 卖空所得可继续交易或归还金圆券借款；消费前还清两种欠款及利息。
  \item 用持有 20 张却卖出 30 张解释新增欠款，合并平多、平空只退出已有持仓的提醒。
\end{itemize}

\section{借款与强制平仓}
\begin{itemize}
  \item 解释杠杆，用 500 本金加 1,000 借款、跌价后剩 110 的例子说明损失。
  \item 正文只解释保证金比率的用途，保留初始门槛与维持门槛的操作对照表。
  \item 借金圆券或欠卡牌都可能强平，无两种欠款时不会仅因价格变化强平。
  \item 说明提前查看平仓报价、准备回补现金，以及强平的费用和剩余债务。
  \item 风险基数、90\% 系数、11.11\% 推导和完整算例移入附录。
\end{itemize}

\section{到价挂单}
\begin{itemize}
  \item 从等低价买入的用途引出挂单，再讲条件、方向、数量和比例。
  \item 保留止盈止损例子；比例按执行时持仓计算，订单独立执行。
  \item 3 张有效挂单、不预留资产、借款授权和额度限制。
  \item 触发不保证成交价格、定时检查可能错过短暂触价、明确失败后结束。
  \item 查记录确认结果；挂单不能保证抢在强平前成交。
\end{itemize}

\section{消费与活动结束}
\begin{itemize}
  \item 兑换与付费弹幕计入累计消费；消费时还清两种欠款及利息，可保留已有卡牌。
  \item 保留消费 200、剩下 450 的例子，明确后续亏损或新增债务不扣减已取得的消费成绩。
  \item 初始本金与额外发放的金圆券可消费计分，费用和未消费资产不计分。
  \item 先停交易，再停消费；卖出和回补须在交易截止前完成。
  \item 不统一按收盘价兑付，不承诺自动平仓、补足本金或免债。
\end{itemize}

\appendix
\section{恒定乘积 AMM 与保证金计算}
\begin{itemize}
  \item 储备、参考价、普通交易的恒定乘积公式；注撤资会改变乘积。
  \item 买卖公式、订单大小算例、手续费和舍入说明。
  \item 保证金比率按金额对照表、还债后净值、分母为何取较大值、门槛比较的顺序讲解；说明避免借款买入重复相加，也不漏算自有持仓或借款；保留 90\% 系数和 11.11\% 门槛推导。
  \item 正文借款例子的完整风险基数和比率计算。
\end{itemize}

\section*{免责声明}
保留原有同人性质、虚拟资产和活动范围等条款。

% 活动规则：名义最高杠杆 10 倍；消费榜按累计消费排名，不扣减当前金圆券债务。
